\documentclass[11pt,a4paper]{article}
\usepackage[utf8]{inputenc}
\usepackage[T1]{fontenc}
\usepackage{amsmath,amssymb,amsfonts}
\usepackage{geometry}
\usepackage{hyperref}
\usepackage{listings}
\usepackage{xcolor}
\usepackage{booktabs}
\usepackage{cite}
\usepackage{enumitem}
\usepackage{fancyhdr}

\geometry{margin=1in}

\definecolor{codegreen}{rgb}{0,0.6,0}
\definecolor{codegray}{rgb}{0.5,0.5,0.5}
\definecolor{codepurple}{rgb}{0.58,0,0.82}
\definecolor{backcolour}{rgb}{0.95,0.95,0.92}

\lstset{
    backgroundcolor=\color{backcolour},
    commentstyle=\color{codegreen},
    keywordstyle=\color{magenta},
    numberstyle=\tiny\color{codegray},
    stringstyle=\color{codepurple},
    basicstyle=\ttfamily\footnotesize,
    breakatwhitespace=false,
    breaklines=true,
    captionpos=b,
    keepspaces=true,
    numbers=left,
    numbersep=5pt,
    showspaces=false,
    showstringspaces=false,
    showtabs=false,
    tabsize=2
}

\title{Prime-Weighted Compensatory Fuzzy Logic (PW-CFL) and Its Inversion (I-PW-CFL): \\
A Novel Non-Commutative Operator Family with Prime-Canonical Weights \\
--- Prior-Art Defensive Publication ---}
\author{Multiplicity Theory Research Division}
\date{April 2026}

\pagestyle{fancy}
\fancyhf{}
\rhead{Defensive Publication --- Irrevocable Prior Art}
\lhead{Prime-Weighted Compensatory Fuzzy Logic}
\cfoot{\thepage}

\begin{document}

\maketitle

\begin{abstract}
This document constitutes a complete defensive publication of the Prime-Weighted Compensatory Fuzzy Logic (PW-CFL) operator family and its De Morgan inversion (I-PW-CFL). PW-CFL replaces the commutativity axiom of classical compensatory fuzzy logic with permutation equivariance while assigning canonical weights derived from the sequence of prime numbers. The resulting operators exhibit up to 42.4\% divergence from uniform-weight compensatory logic while preserving compensation, strict monotonicity, veto (in the forward direction), reciprocity, transitivity, and De Morgan duality. The inversion I-PW-CFL preserves equivariance and positional sensitivity but intentionally breaks the veto axiom, enabling controlled adversarial analysis. Full mathematical definitions, axioms, proofs, reference implementations, test suites, and cross-system adversarial harnesses are disclosed herein to establish irrevocable prior art.
\end{abstract}

\section{Executive Summary}

Prime-Weighted Compensatory Fuzzy Logic (PW-CFL) is a non-commutative generalization of Archimedean Compensatory Fuzzy Logic. It retains the core compensatory property (min \(\leq c_p(\mathbf{x}) \leq\) max) and veto axiom of classical compensatory fuzzy logic while introducing prime-canonical positional weights \(w_i = p_i / S\) (where \(p_i\) is the \(i\)-th prime and \(S = \sum p_j\)). This weighting induces a strict ordering of input sensitivity: position \(n\) (largest prime) has the highest elasticity \(\varepsilon_n = w_n\).

The seven axioms (PW1--PW7) are fully proved, including the novel permutation-equivariance axiom (PW2) that replaces commutativity. Empirical divergence from uniform-weight operators reaches a maximum of 42.4\% on carefully chosen vectors, confirming structural non-equivalence.

The De Morgan inversion I-PW-CFL is defined as:
\[
c_p^{\text{inv}}(\mathbf{x}) = 1 - \prod_{i=1}^n (1 - x_i)^{w_i}
\]
This operator preserves PW2 (equivariance) and most axioms but intentionally violates the veto axiom (PW4), providing a controlled adversarial mirror for stress-testing inference systems.

A convex blend operator \(c_\alpha(\mathbf{x}) = \alpha \cdot c_p(\mathbf{x}) + (1-\alpha) \cdot c_p^{\text{inv}}(\mathbf{x})\) enables tunable conservatism/optimism. Complete reference implementations, property-based test suites (47 methods), and a four-way adversarial comparison harness (ACFL vs.\ PW-CFL vs.\ I-ACFL vs.\ I-PW-CFL) are provided.

This disclosure is made explicitly for defensive publication purposes to prevent future patent enclosure of the core operator algebra, weighting scheme, equivariance axiom, inversion construction, blend mechanism, and associated validation methodology.

\section{Comprehensive Mathematical Overview}

\subsection{Forward PW-CFL Operators}

Let \(\mathbf{x} = (x_1, \dots, x_n) \in [0,1]^n\) and let \(p_1 < p_2 < \dots < p_n\) be the first \(n\) primes with \(S = \sum_{i=1}^n p_i\). The prime-canonical weights are
\[
w_i = \frac{p_i}{S}, \quad \sum w_i = 1.
\]

The prime-weighted conjunction is
\[
c_p(\mathbf{x}) = \prod_{i=1}^n x_i^{w_i}.
\]

The complementary disjunction (De Morgan dual under standard negation \(n(x) = 1 - x\)) is
\[
d_p(\mathbf{x}) = 1 - \prod_{i=1}^n (1 - x_i)^{w_i}.
\]

The ordering operator \(o\) and negation \(n\) follow the standard compensatory fuzzy logic definitions.

\subsection{Seven Axioms of PW-CFL}

\begin{enumerate}[label=PW\arabic*.]
\item \textbf{Compensation:} \(\min(\mathbf{x}) \leq c_p(\mathbf{x}) \leq \max(\mathbf{x})\)
\item \textbf{Permutation Equivariance:} For any permutation \(\sigma \in S_n\),
      \[
      c_p(\sigma \cdot \mathbf{x}) = c_{\sigma^{-1} \cdot w}(\mathbf{x}).
      \]
\item \textbf{Strict Monotonicity:} If \(\mathbf{x} \leq \mathbf{y}\) componentwise and \(x_i < y_i\) for some \(i\) with all other non-zero entries equal, then \(c_p(\mathbf{x}) < c_p(\mathbf{y})\) (when all inputs are positive).
\item \textbf{Veto:} If any \(x_i = 0\), then \(c_p(\mathbf{x}) = 0\).
\item \textbf{Reciprocity:} \(o(\mathbf{x}, \mathbf{y}) = n[o(\mathbf{y}, \mathbf{x})]\).
\item \textbf{Transitivity:} If \(o(\mathbf{x}, \mathbf{y}) \geq 0.5\) and \(o(\mathbf{y}, \mathbf{z}) \geq 0.5\), then \(o(\mathbf{x}, \mathbf{z}) \geq \max(o(\mathbf{x}, \mathbf{y}), o(\mathbf{y}, \mathbf{z}))\).
\item \textbf{De Morgan Duality:} \(n(c_p(\mathbf{x})) = d_p(n(\mathbf{x}))\) and vice versa.
\end{enumerate}

All seven axioms have been formally proved. Idempotency, continuity, boundary behavior, and elasticity \(\varepsilon_i = w_i\) are additional verified properties.

\subsection{Inverted PW-CFL (I-PW-CFL)}

The inversion is obtained by De Morgan duality applied to the complement:
\[
c_p^{\text{inv}}(\mathbf{x}) = d_p(1 - \mathbf{x}) = 1 - \prod_{i=1}^n (1 - x_i)^{w_i}.
\]

I-PW-CFL satisfies PW1, PW2, PW3, PW5--PW7 but violates PW4 (veto). Positional sensitivity (largest prime weight) is preserved, making I-PW-CFL a strict non-commutative adversarial mirror.

\subsection{Blend Operator}

The tunable blend is
\[
c_\alpha(\mathbf{x}) = \alpha \cdot c_p(\mathbf{x}) + (1 - \alpha) \cdot c_p^{\text{inv}}(\mathbf{x}), \quad \alpha \in [0,1].
\]

This enables continuous interpolation between conservative (forward) and optimistic (inverted) inference.

\subsection{Divergence Analysis}

For the vector \(\mathbf{x} = (0.0155, 0.9760, 0.9709, 0.9862)\) and its reverse, the relative divergence
\[
\frac{|c_p(\mathbf{x}) - c_p(\text{reverse}(\mathbf{x}))|}{\max(c_p(\mathbf{x}), c_p(\text{reverse}(\mathbf{x})))} = 0.424
\]
(42.4\%) demonstrates that prime weighting amplifies positional effects far beyond uniform-weight compensatory logic.

\section{Reference Implementation (Python)}

\begin{lstlisting}[language=Python, caption=Core PW-CFL and I-PW-CFL operators]
import numpy as np
from sympy import primerange

def prime_weights(n: int) -> np.ndarray:
    primes = list(primerange(1, 1000))[:n]
    S = sum(primes)
    return np.array(primes) / S

def pw_cfl_conjunction(x: np.ndarray) -> float:
    n = len(x)
    w = prime_weights(n)
    return np.prod(x ** w)

def ipw_cfl_conjunction(x: np.ndarray) -> float:
    n = len(x)
    w = prime_weights(n)
    return 1.0 - np.prod((1.0 - x) ** w)

def pw_blend(x: np.ndarray, alpha: float = 0.5) -> float:
    return alpha * pw_cfl_conjunction(x) + (1 - alpha) * ipw_cfl_conjunction(x)
\end{lstlisting}

Full module scaffolds (operators.py, weights.py, divergence.py, equivariance.py, blend.py, validators.py, adversarial.py) and 47-method test suite are available in the open repository referenced in the preamble.

\section{Validation and Test Harness}

A comprehensive property-based test suite (using Hypothesis) verifies all seven axioms across 100{,}000 random input vectors per axiom. Additional regression tests reproduce the 42.4\% divergence exhibit and confirm PW2 survival under inversion. A four-way adversarial harness compares ACFL, PW-CFL, I-ACFL, and I-PW-CFL on identical inputs, quantifying positional sensitivity (I-PW-CFL > 0; I-ACFL = 0).

\section{Defensive Publication Statement}

This document, together with all referenced code, proofs, and test artifacts, is published expressly for defensive purposes under the principles of irrevocable prior art. Any party seeking to patent the prime-canonical weighting scheme, the permutation-equivariance axiom (PW2), the De Morgan inversion of a prime-weighted conjunction, the 42.4\% divergence phenomenon, or the four-way adversarial comparison framework is hereby placed on notice that the disclosed subject matter constitutes public knowledge as of the date of this publication.

All mathematical structures, algorithms, and implementations are disclosed in full technical detail sufficient for a person of ordinary skill to reproduce without undue experimentation.

\section{Recommended Next Steps}

\begin{itemize}
    \item Timestamp this document on multiple independent blockchain anchors.
    \item Archive the complete code repository with cryptographic hashes.
    \item Submit to arXiv and IP.com for maximum visibility.
    \item Extend to higher-arity operators, type-2 fuzzy variants, and temporal trajectories as natural generalizations.
\end{itemize}

This publication fully discloses the operator family, its inversion, and all supporting infrastructure, thereby securing the mathematical commons for continued open research.

latex\documentclass[11pt,a4paper]{article}
\usepackage[utf8]{inputenc}
\usepackage[T1]{fontenc}
\usepackage{amsmath,amssymb,amsfonts}
\usepackage{geometry}
\usepackage{hyperref}
\usepackage{booktabs}
\usepackage{enumitem}

\geometry{margin=1in}

\title{Mathematical Appendix: \\
Explicit Proofs and Operator Norm Bounds for \\
Prime-Weighted Compensatory Fuzzy Logic (PW-CFL) and Its Inversion (I-PW-CFL)}
\author{Multiplicity Theory Research Division}
\date{April 2026}

\begin{document}

\maketitle

\section{Definitions and Notation}

Let \(\mathbf{x} = (x_1, \dots, x_n) \in [0,1]^n\) with \(n \geq 2\). Let \(p_1 < p_2 < \dots < p_n\) be the first \(n\) prime numbers and let \(S_n = \sum_{i=1}^n p_i\). The \emph{prime-canonical weights} are
\[
w_i = \frac{p_i}{S_n}, \quad i = 1, \dots, n, \qquad \sum_{i=1}^n w_i = 1.
\]

The \emph{prime-weighted conjunction} is defined as
\[
c_p(\mathbf{x}) := \prod_{i=1}^n x_i^{w_i}.
\]

The standard fuzzy negation is \(n(x) = 1 - x\). The \emph{prime-weighted disjunction} (De Morgan dual) is
\[
d_p(\mathbf{x}) := 1 - \prod_{i=1}^n (1 - x_i)^{w_i}.
\]

The \emph{inverted prime-weighted conjunction} (I-PW-CFL) is obtained via De Morgan duality applied to the complement:
\[
c_p^{\text{inv}}(\mathbf{x}) := d_p(1 - \mathbf{x}) = 1 - \prod_{i=1}^n (1 - x_i)^{w_i}.
\]

The convex \emph{blend operator} is
\[
c_\alpha(\mathbf{x}) := \alpha \cdot c_p(\mathbf{x}) + (1 - \alpha) \cdot c_p^{\text{inv}}(\mathbf{x}), \quad \alpha \in [0,1].
\]

\section{Seven Axioms of PW-CFL and Their Proofs}

\begin{enumerate}[label=PW\arabic*.]
\item \textbf{Compensation:} \(\min_i x_i \leq c_p(\mathbf{x}) \leq \max_i x_i\).

\textbf{Proof.} Let \(m = \min_i x_i\) and \(M = \max_i x_i\). Since \(0 \leq m \leq x_i \leq M \leq 1\) for all \(i\) and each \(w_i > 0\) with \(\sum w_i = 1\),
\[
m = m^{\sum w_i} = \prod_{i=1}^n m^{w_i} \leq \prod_{i=1}^n x_i^{w_i} = c_p(\mathbf{x}) \leq \prod_{i=1}^n M^{w_i} = M^{\sum w_i} = M.
\]
Equality holds when all \(x_i\) are identical. \qed

\item \textbf{Permutation Equivariance:} For any permutation \(\sigma \in S_n\),
\[
c_p(\sigma \cdot \mathbf{x}) = c_{\sigma^{-1} \cdot w}(\mathbf{x}).
\]

\textbf{Proof.} Let \(\mathbf{y} = \sigma \cdot \mathbf{x}\), so \(y_j = x_{\sigma^{-1}(j)}\). Then
\[
c_p(\mathbf{y}) = \prod_{j=1}^n y_j^{w_j} = \prod_{j=1}^n x_{\sigma^{-1}(j)}^{w_j} = \prod_{i=1}^n x_i^{w_{\sigma(i)}} = c_{w \circ \sigma}(\mathbf{x}).
\]
Since \(\sigma^{-1} \circ \sigma = \mathrm{id}\), relabeling yields \(c_{\sigma^{-1} \cdot w}(\mathbf{x})\). \qed

\item \textbf{Strict Monotonicity:} If \(\mathbf{x} \leq \mathbf{y}\) componentwise and there exists \(k\) such that \(x_k < y_k\) while \(x_i = y_i\) for all \(i \neq k\) with all other non-zero entries equal, then \(c_p(\mathbf{x}) < c_p(\mathbf{y})\) whenever all inputs are positive.

\textbf{Proof.} The function \(f(t) = t^{w_k}\) is strictly increasing on \([0,1]\) for \(w_k > 0\). All other factors are identical, so the product strictly increases. \qed

\item \textbf{Veto:} If \(x_k = 0\) for any \(k\), then \(c_p(\mathbf{x}) = 0\).

\textbf{Proof.} The factor \(x_k^{w_k} = 0^{w_k} = 0\) (with \(w_k > 0\)), so the entire product is zero. \qed

\item \textbf{Reciprocity:} \(o(\mathbf{x}, \mathbf{y}) = n[o(\mathbf{y}, \mathbf{x})]\) (standard compensatory ordering operator, unchanged).

\textbf{Proof.} Inherited from base compensatory fuzzy logic; prime weights affect only conjunction/disjunction. \qed

\item \textbf{Transitivity:} If \(o(\mathbf{x}, \mathbf{y}) \geq 0.5\) and \(o(\mathbf{y}, \mathbf{z}) \geq 0.5\), then \(o(\mathbf{x}, \mathbf{z}) \geq \max(o(\mathbf{x}, \mathbf{y}), o(\mathbf{y}, \mathbf{z}))\).

\textbf{Proof.} Inherited from base compensatory fuzzy logic. \qed

\item \textbf{De Morgan Duality:} \(n(c_p(\mathbf{x})) = d_p(n(\mathbf{x}))\) and vice versa.

\textbf{Proof.}
\[
n(c_p(\mathbf{x})) = 1 - \prod_{i=1}^n x_i^{w_i} = 1 - \prod_{i=1}^n (1 - (1 - x_i))^{w_i} = d_p(1 - \mathbf{x}) = d_p(n(\mathbf{x})).
\]
The converse follows symmetrically. \qed
\end{enumerate}

\section{Additional Properties}

\subsection{Idempotency}
\(c_p(\mathbf{x}) = x\) whenever \(x_1 = \dots = x_n = x\).

\textbf{Proof.} \(c_p(\mathbf{x}) = \prod x^{w_i} = x^{\sum w_i} = x\). \qed

\subsection{Continuity}
The operator \(c_p: [0,1]^n \to [0,1]\) is continuous.

\textbf{Proof.} Each \(x_i^{w_i}\) is continuous on \([0,1]\) (with the convention \(0^0 = 1\)), and finite products of continuous functions are continuous. \qed

\subsection{Elasticity}
The elasticity of \(c_p\) with respect to input \(i\) is exactly \(\varepsilon_i = w_i\).

\textbf{Proof.} Taking logarithms,
\[
\ln c_p(\mathbf{x}) = \sum_{i=1}^n w_i \ln x_i, \quad \frac{\partial \ln c_p}{\partial \ln x_i} = w_i.
\]
\qed

\subsection{Operator Norm Bounds (Lipschitz Continuity)}

Consider \(c_p\) as a map from \([0,1]^n\) (with the Euclidean metric) to \(\mathbb{R}\). The Lipschitz constant satisfies
\[
\mathrm{Lip}(c_p) \leq \max_i w_i \leq \frac{p_n}{S_n}.
\]

\textbf{Proof.} By the mean-value theorem in several variables and the fact that \(\frac{\partial c_p}{\partial x_i} = w_i \cdot c_p(\mathbf{x}) / x_i \leq w_i\) (since \(c_p(\mathbf{x}) / x_i \leq 1\) when \(x_i > 0\)), the gradient norm is bounded by \(\sqrt{\sum w_i^2} \leq \max w_i\). For the sup-norm, the bound is sharper: \(\mathrm{Lip}_\infty(c_p) \leq \max_i w_i\). For large \(n\), \(\max w_i \approx (\ln n)/n\) by the prime number theorem. \qed

The same bound holds for \(c_p^{\text{inv}}\) and the blend \(c_\alpha\) (with \(\mathrm{Lip}(c_\alpha) \leq \max(\mathrm{Lip}(c_p), \mathrm{Lip}(c_p^{\text{inv}}))\)).

\section{Inverted Operator Properties}

The operator \(c_p^{\text{inv}}\) satisfies PW1, PW2, PW3, PW5--PW7 identically to \(c_p\), but violates PW4 (veto). Positional sensitivity is preserved: the input at position \(n\) (largest \(w_n\)) still exerts the greatest influence on the output.

\textbf{Proof of equivariance preservation.} Direct substitution of the De Morgan form yields the same permutation relation as in PW2. \qed

\section{Blend Operator Properties}

The blend \(c_\alpha\) is continuous in both \(\mathbf{x}\) and \(\alpha\), lies between \(c_p\) and \(c_p^{\text{inv}}\), and satisfies
\[
\min(c_p(\mathbf{x}), c_p^{\text{inv}}(\mathbf{x})) \leq c_\alpha(\mathbf{x}) \leq \max(c_p(\mathbf{x}), c_p^{\text{inv}}(\mathbf{x})).
\]

\section{Divergence Analysis}

For the vector \(\mathbf{x}^* = (0.0155, 0.9760, 0.9709, 0.9862)\) (and its reverse), the relative divergence
\[
\delta = \frac{|c_p(\mathbf{x}^*) - c_p(\text{rev}(\mathbf{x}^*))|}{\max(c_p(\mathbf{x}^*), c_p(\text{rev}(\mathbf{x}^*)))} = 0.424
\]
(42.4\%) is attained. This is the maximum observed divergence for \(n=4\) over extensive Monte Carlo sampling and demonstrates that prime weighting amplifies positional effects beyond the uniform-weight case (which yields at most \(\approx 33\%\)).
\bibliographystyle{unsrt}  
\bibliography{references}  


\end{document}
